\documentclass{article}
\usepackage[utf8]{inputenc}
\usepackage{amsmath, amssymb}
\usepackage{graphicx}
\usepackage{hyperref}
\usepackage{siunitx}
\usepackage{booktabs}
\title{Teori for håndholdt EMP-generator}
\author{}
\date{}

\begin{document}
\maketitle

\section{Innledning}
En elektromagnetisk puls (EMP) er en kort, intens utbrudd av elektromagnetisk energi. I dette prosjektet skapes pulsen ved å utlade en stor kondensator gjennom en spole med svært lav induktans. Den raske endringen i strøm ($di/dt$) genererer et tidsvarierende magnetfelt, som ifølge Faradays lov induserer en spenning i nærliggende ledere. Hvis den induserte spenningen overstiger tålegrensen til elektroniske komponenter, oppstår permanent skade.

\section{Grunnleggende fysikk}

\subsection{Energilagring i kondensator}
Energien som er lagret i en kondensator med kapasitans $C$ og spenning $V$ er gitt ved
\begin{equation}
E = \frac{1}{2} C V^2.
\end{equation}
For $C = 1000\,\mu\text{F}$ og $V = 400\,\text{V}$ får vi
\begin{equation}
E = 0.5 \cdot 1000 \times 10^{-6} \cdot (400)^2 = 80\,\text{J}.
\end{equation}
80 joule tilsvarer energien i et kraftig luftgevær, men her leveres den på mikrosekunder.

\subsection{Utladning gjennom spole}
Når kondensatoren kobles over en spole med induktans $L$ og ohmsk motstand $R$, dannes en serie RLC-krets. Strømmen $i(t)$ er gitt av differensialligningen
\begin{equation}
L \frac{d^2 i}{dt^2} + R \frac{di}{dt} + \frac{1}{C} i = 0.
\end{equation}
For en underdempet krets ($R < 2\sqrt{L/C}$) er løsningen en dempet sinus:
\begin{equation}
i(t) = \frac{V_0}{\omega L} e^{-\alpha t} \sin(\omega t),
\end{equation}
der
\begin{align}
\alpha &= \frac{R}{2L}, \\
\omega &= \sqrt{\frac{1}{LC} - \alpha^2}.
\end{align}
Toppstrømmen inntreffer ved $t \approx \pi/(2\omega)$ og er tilnærmet
\begin{equation}
I_{\text{peak}} \approx V_0 \sqrt{\frac{C}{L}}.
\end{equation}
Med $L = 1.9\,\mu\text{H}$ (se \texttt{beregninger.tex}) og $C = 1000\,\mu\text{F}$ får vi
\begin{equation}
I_{\text{peak}} \approx 400 \sqrt{\frac{1000\times10^{-6}}{1.9\times10^{-6}}} = 400 \times \sqrt{526.3} \approx 9\,174\,\text{A}.
\end{equation}
I praksis vil parasittisk motstand (ESR i kondensator, ledningsmotstand, lysbuemotstand) redusere toppstrømmen til ca. $8\,000$--$9\,000$ A. Disse verdiene er teoretiske estimater og ikke eksperimentelt bekreftet.

\subsection{Magnetfelt fra en flat spiralspole}
For en flat, sirkulær spole med $N$ vindinger, middelradius $r$, og strøm $I$, er magnetfeltet i sentrum tilnærmet
\begin{equation}
B = \frac{\mu_0 N I}{2r},
\end{equation}
der $\mu_0 = 4\pi \times 10^{-7}\,\text{H/m}$ er permeabiliteten i vakuum. Med $N=6$, $r=5\,\text{cm}$ og $I=9\,000\,\text{A}$:
\begin{equation}
B = \frac{4\pi \times 10^{-7} \cdot 6 \cdot 9000}{2 \cdot 0.05} \approx 0.679\,\text{T}.
\end{equation}
Til sammenligning er jordens magnetfelt ca. $50\,\mu\text{T}$, så dette feltet er over 13 000 ganger sterkere.

\subsection{Indusert spenning i en offerkrets}
En liten lukket sløyfe med areal $A$ som gjennomtrenges av et tidsvarierende magnetfelt $B(t)$, får indusert en elektromotorisk spenning (ems) gitt av Faradays lov:
\begin{equation}
\mathcal{E} = -\frac{d\Phi}{dt} = -A \frac{dB}{dt}.
\end{equation}
Hvis feltet endrer seg fra 0 til $B_{\text{max}}$ på tiden $\Delta t$, blir den gjennomsnittlige induserte spenningen
\begin{equation}
|\mathcal{E}| \approx A \frac{B_{\text{max}}}{\Delta t}.
\end{equation}
For en sløyfe med $A = 1\,\text{cm}^2 = 1\times10^{-4}\,\text{m}^2$, $B_{\text{max}} = 0.06\,\text{T}$ (ved 10 cm avstand, estimert) og $\Delta t = 68.5\,\mu\text{s}$ (kvart periode):
\begin{equation}
|\mathcal{E}| \approx 1\times10^{-4} \cdot \frac{0.06}{68.5\times10^{-6}} \approx 0.088\,\text{V}.
\end{equation}
Dette er for en enkelt, liten sløyfe. I virkelige kretser med lengre ledningsbaner, flere sløyfer og antenneeffekter kan spenningen bli mange titalls eller hundrevis av volt – nok til å ødelegge halvledere. Merk at den teoretiske stigetiden kan være kortere i praksis, noe som gir høyere indusert spenning.

\section{Kretsens virkemåte}

\subsection{Overordnet blokkskjema}
Enheten består av tre hoveddeler:
\begin{enumerate}
\item \textbf{Ladekrets:} En flyback-omformer som omdanner 54 V DC fra batteriene til 400 V DC for å lade hovedkondensatoren.
\item \textbf{Energilager og utladningskrets:} En 1000 µF kondensator og et trigget gnistgap som kobler kondensatoren til spolen.
\item \textbf{Antennespole:} En flat spiralspole som omdanner strømpulsen til et magnetfelt.
\end{enumerate}

\subsection{Ladekrets (flyback-omformer)}
Flyback-omformeren er en selvoscillerende krets basert på en MOSFET, en ferritttransformator og en diode. Når batterispenningen påtrykkes, begynner MOSFET-en å lede på grunn av lekkasjestrøm. Strømmen gjennom primærviklingen induserer en spenning i tilbakekoblingsviklingen som driver gaten positivt, og MOSFET-en åpner helt. Strømmen i primærviklingen øker lineært, og energi lagres i transformatorens magnetfelt. Når kjernen begynner å mettes, kollapser tilbakekoblingen, MOSFET-en slår seg av, og den lagrede energien overføres til sekundærviklingen. Utgangsspenningen likerettes av en rask diode og lader hovedkondensatoren. Prosessen gjentas så lenge ladeknappen holdes inne.

\subsection{Gnistgap og trigger}
Et gnistgap består av to elektroder med en liten luftavstand. Når spenningen over gapet overstiger gjennomslagsspenningen (ca. 3 kV/mm for tørr luft), ioniseres luften og en gnist dannes. For å kontrollere nøyaktig når gnisten utløses, brukes et trigget gnistgap: en tredje elektrode (trigger) plasseres mellom hovedelektrodene. En piezoelektrisk tenner (fra en engangsgrill) genererer en høy spenningspuls (>10 kV) som lager en liten gnist fra triggeren til en av hovedelektrodene. Dette ioniserer gapet, slik at hovedgnisten kan dannes selv om spenningen er lavere enn den naturlige gjennomslagsspenningen.

\subsection{Utladningsforløp}
Når gnistgapet lukkes, kobles den oppladede kondensatoren direkte over spolen. Kretsen er en underdempet RLC-krets med svært lav motstand (dominert av gnistgapets lysbue og ledningsmotstand). Strømmen stiger til over 9 kA i løpet av noen titalls mikrosekunder, og magnetfeltet når sitt maksimum. Deretter svinger energien mellom kondensatoren og spolen, og pulsen dør ut etter noen få perioder på grunn av resistive tap.

\section{Komponentvalg og begrunnelser}
Se \texttt{design\_decisions.tex} for en detaljert diskusjon.
\end{document}
